\section*{第 \ref{ch:g7:circuits-series-parallel} 章——电路：回路、串联与并联}

\begin{solution}{exo:g7:circuits-series-parallel:1}
电荷在\omterm{def:g3:first-electric-circuit:battery}{电池}的推动下沿电路定向流动，
这就是\omterm{def:g7:circuits-series-parallel:current}{电流}。只有当至少一个闭合回路
连接两个\omterm{def:g3:first-electric-circuit:battery}{接线柱}时，\omterm{def:g7:circuits-series-parallel:current}{电流}才存在；断开回路，
\omterm{def:g7:circuits-series-parallel:current}{电流}便同时在各处停止。
\end{solution}

\begin{solution}{exo:g7:circuits-series-parallel:2}
在\omterm{def:g3:first-electric-circuit:battery}{电池}外部，\omterm{def:g7:circuits-series-parallel:current}{电流}由 $+$ 极出发，沿电路
流回 $-$ 极。在只有一个\omterm{def:g3:first-electric-circuit:bulb}{灯泡}的回路中，箭头
从 $+$ 极出发，经过\omterm{def:g3:first-electric-circuit:bulb}{灯泡}，回到 $-$ 极。
\end{solution}

\begin{solution}{exo:g7:circuits-series-parallel:3}
一种常见的错误想法认为用电器会消耗\omterm{def:g7:circuits-series-parallel:current}{电流}，
因此“下游”的\omterm{def:g3:first-electric-circuit:bulb}{灯泡}应当更暗。亮度相同——
无论怎样对调都如此——说明两个\omterm{def:g3:first-electric-circuit:bulb}{灯泡}通过的是
同一股完整的\omterm{def:g7:circuits-series-parallel:current}{电流}，途中没有\omterm{def:g7:circuits-series-parallel:current}{电流}被消耗。
\end{solution}

\begin{solution}{exo:g7:circuits-series-parallel:4}
两者是同一股\omterm{def:g7:circuits-series-parallel:current}{电流}：
单一回路中，一股\omterm{def:g7:circuits-series-parallel:current}{电流}流经其上的每个用电器。
\end{solution}

\begin{solution}{exo:g7:circuits-series-parallel:5}
\omterm{def:g7:circuits-series-parallel:current}{电流}在\omterm{def:g5:series-parallel-circuits:parallel}{节点}处分开——一部分走每条\omterm{def:g5:series-parallel-circuits:parallel}{支路}，
在第二个\omterm{def:g5:series-parallel-circuits:parallel}{节点}处各部分重新汇合，恢复完整。
每条\omterm{def:g5:series-parallel-circuits:parallel}{支路}都获得\omterm{def:g3:first-electric-circuit:battery}{电池}的全部推动。
\end{solution}

\begin{solution}{exo:g7:circuits-series-parallel:6}
单独 $=$ 与一个\omterm{def:g5:series-parallel-circuits:parallel}{并联}邻居（亮度最大），
再\omterm{def:g4:series-circuits:series}{串联}一个伙伴（变暗），
然后\omterm{def:g4:series-circuits:series}{串联}两个伙伴（更暗）。
\end{solution}

\begin{solution}{exo:g7:circuits-series-parallel:7}
\omterm{def:g5:series-parallel-circuits:parallel}{并联}：两条\omterm{def:g5:series-parallel-circuits:parallel}{支路}各自通过一股完整的\omterm{def:g7:circuits-series-parallel:current}{电流}，
使\omterm{def:g3:first-electric-circuit:battery}{电池}的总输出加倍；而\omterm{def:g4:series-circuits:series}{串联}的一对\omterm{def:g3:first-electric-circuit:bulb}{灯泡}
只让一股受限的\omterm{def:g7:circuits-series-parallel:current}{电流}依次通过两者。
满亮度的代价由\omterm{def:g3:first-electric-circuit:battery}{电池}承担。
\end{solution}

\begin{solution}{exo:g7:circuits-series-parallel:8}
缆车上的椅子代表电荷，它们周而复始地循环，
从不被消耗；滑雪者代表\omterm{def:g5:everyday-energy:energy}{能量}，在\omterm{def:g3:first-electric-circuit:battery}{电池}处上车，
在\omterm{def:g3:first-electric-circuit:bulb}{灯泡}处以光和热的形式下车。耗尽的是
滑雪场中滑雪者的存量，即\omterm{def:g3:first-electric-circuit:battery}{电池}所储存的\omterm{def:g5:everyday-energy:energy}{能量}。
\end{solution}

\begin{solution}{exo:g7:circuits-series-parallel:9}
经过 A 的每个回路都先经过\omterm{ex:g3:first-electric-circuit:switch}{开关}，再经过 B 或 C——
因此 A 的\omterm{def:g4:series-circuits:series}{串联}伙伴是\omterm{ex:g3:first-electric-circuit:switch}{开关}和这\emph{一对}\omterm{def:g5:series-parallel-circuits:parallel}{支路}整体，
而不是单独的 B 或 C。断开 C 所在的\omterm{def:g5:series-parallel-circuits:parallel}{支路}，
整股\omterm{def:g7:circuits-series-parallel:current}{电流}便全部流经 B，A 仍然发光，
只是亮度略有不同——只剩一条通路，而非两条——
这一点将由明年的仪器测出。
\end{solution}

\begin{solution}{exo:g7:circuits-series-parallel:10}
电源之后的公共干路上设一个总\omterm{ex:g3:first-electric-circuit:switch}{开关}，
随后是十条\omterm{def:g5:series-parallel-circuits:parallel}{并联}\omterm{def:g5:series-parallel-circuits:parallel}{支路}，每条\omterm{def:g5:series-parallel-circuits:parallel}{支路}接一只\omterm{def:g3:first-electric-circuit:bulb}{灯泡}。
\omterm{def:g5:series-parallel-circuits:parallel}{并联}规律使每只\omterm{def:g3:first-electric-circuit:bulb}{灯泡}都获得完全独立的亮度；
总\omterm{ex:g3:first-electric-circuit:switch}{开关}串接在所有回路共有的干路上，
因此一个\omterm{ex:g3:first-electric-circuit:switch}{开关}即可同时控制全部十只\omterm{def:g3:first-electric-circuit:bulb}{灯泡}。
\end{solution}

\begin{solution}{exo:g7:circuits-series-parallel:11}
\omterm{def:g5:series-parallel-circuits:parallel}{节点}处的分流并不是均分：每条\omterm{def:g5:series-parallel-circuits:parallel}{支路}自身的
通路可能阻碍\omterm{def:g7:circuits-series-parallel:current}{电流}，过长过细的导线会限制
该\omterm{def:g5:series-parallel-circuits:parallel}{支路}分得的份额——每秒通过的电荷更少，
灯光也更弱。通路本身与\omterm{def:g7:circuits-series-parallel:current}{电流}相抗衡，
这便是明年要学的电阻概念的基础。
\end{solution}

\begin{solution}{exo:g7:circuits-series-parallel:12}
X 独占一条\omterm{def:g5:series-parallel-circuits:parallel}{支路}，获得全部推动，最亮。Y 与 Z
在同一条\omterm{def:g5:series-parallel-circuits:parallel}{支路}中\omterm{def:g4:series-circuits:series}{串联}，分享推动，亮度相同，
且都比 X 暗。若 Z 烧坏，其回路断开，Y 随之变暗，
而 X 依旧发光，毫不受影响。若 X 烧坏，只有 X 所在\omterm{def:g5:series-parallel-circuits:parallel}{支路}
失效，Y 与 Z 照常发光，亮度与先前相同。
\end{solution}

\begin{solution}{pb:g7:circuits-series-parallel:1}
\textbf{1.} 四条\omterm{def:g5:series-parallel-circuits:parallel}{并联}\omterm{def:g5:series-parallel-circuits:parallel}{支路}，各负责一项任务：\omterm{def:g5:series-parallel-circuits:parallel}{并联}时，
每条\omterm{def:g5:series-parallel-circuits:parallel}{支路}无论其他\omterm{def:g5:series-parallel-circuits:parallel}{支路}如何，都以满功率工作。
任何\omterm{def:g4:series-circuits:series}{串联}方式都会让阅读灯的\omterm{ex:g3:first-electric-circuit:switch}{开关}或一只烧坏的
楼梯\omterm{def:g3:first-electric-circuit:bulb}{灯泡}使大灯变暗，甚至熄灭——这是绝不允许的。
\textbf{2.} 总\omterm{ex:g3:first-electric-circuit:switch}{开关}位于\omterm{def:g3:first-electric-circuit:battery}{电池}组之后的公共干路上；
四条\omterm{def:g5:series-parallel-circuits:parallel}{支路}各自的\omterm{ex:g3:first-electric-circuit:switch}{开关}与其用电器\omterm{def:g4:series-circuits:series}{串联}，
位于各自\omterm{def:g5:series-parallel-circuits:parallel}{支路}之内。
\textbf{3.} 是的——依\omterm{def:g5:series-parallel-circuits:parallel}{并联}规律，一条\omterm{def:g5:series-parallel-circuits:parallel}{支路}出现断路，
只会切断该\omterm{def:g5:series-parallel-circuits:parallel}{支路}所分得的\omterm{def:g7:circuits-series-parallel:current}{电流}，其余\omterm{def:g5:series-parallel-circuits:parallel}{支路}照常工作。
\textbf{4.} \omterm{def:g3:first-electric-circuit:battery}{电池}组、总\omterm{ex:g3:first-electric-circuit:switch}{开关}，再经\omterm{def:g5:series-parallel-circuits:parallel}{节点}分为四条
\omterm{def:g5:series-parallel-circuits:parallel}{支路}——（大灯 $+$ \omterm{ex:g3:first-electric-circuit:switch}{开关}）、（雾笛 $+$ \omterm{ex:g3:first-electric-circuit:switch}{开关}）、
（阅读灯 $+$ \omterm{ex:g3:first-electric-circuit:switch}{开关}）、（楼梯灯 $+$ \omterm{ex:g3:first-electric-circuit:switch}{开关}）——
再汇合并回到\omterm{def:g3:first-electric-circuit:battery}{电池}组。
\textbf{5.} 是的：每条\omterm{def:g5:series-parallel-circuits:parallel}{支路}都跨接在\omterm{def:g3:first-electric-circuit:battery}{电池}组两极之间，
因此每条\omterm{def:g5:series-parallel-circuits:parallel}{支路}都获得完整的推动，
不论旁边还有多少\omterm{def:g5:series-parallel-circuits:parallel}{支路}同时取电。
\textbf{6.} 不对：流过\omterm{def:g3:first-electric-circuit:bulb}{灯泡}的\omterm{def:g7:circuits-series-parallel:current}{电流}离开\omterm{def:g3:first-electric-circuit:bulb}{灯泡}时仍然完整——
那些循环往复的“椅子”并未被消耗。\omterm{def:g3:first-electric-circuit:bulb}{灯泡}所取用的是
\omterm{def:g3:first-electric-circuit:battery}{电池}组所储存的\omterm{def:g5:everyday-energy:energy}{能量}，由流过的\omterm{def:g7:circuits-series-parallel:current}{电流}输送而来。
\textbf{7.} 四股完整\omterm{def:g7:circuits-series-parallel:current}{电流}同时流动：\omterm{def:g3:first-electric-circuit:battery}{电池}组的\omterm{def:g5:everyday-energy:energy}{能量}
迅速被四项任务耗尽。代价由\omterm{def:g3:first-electric-circuit:battery}{电池}组承担——
体现为所储存的\omterm{def:g5:everyday-energy:energy}{能量}，而绝非\omterm{def:g7:circuits-series-parallel:current}{电流}的损失。
\textbf{8.} 从\omterm{def:g3:first-electric-circuit:battery}{电池}组的 $+$ 极出发，经过总\omterm{ex:g3:first-electric-circuit:switch}{开关}，
到达\omterm{def:g5:series-parallel-circuits:parallel}{节点}；进入阅读灯所在\omterm{def:g5:series-parallel-circuits:parallel}{支路}，经过其闭合的
\omterm{ex:g3:first-electric-circuit:switch}{开关}和\omterm{def:g3:first-electric-circuit:bulb}{灯泡}；经过\omterm{def:g5:series-parallel-circuits:parallel}{节点}，其余三条（不亮的）
\omterm{def:g5:series-parallel-circuits:parallel}{支路}在此汇合；最后回到\omterm{def:g3:first-electric-circuit:battery}{电池}组的 $-$ 极。
它完全避开了大灯、雾笛和楼梯灯所在的\omterm{def:g5:series-parallel-circuits:parallel}{支路}。
\textbf{9.} 别无其他——只有楼梯灯\omterm{def:g5:series-parallel-circuits:parallel}{支路}随其\omterm{def:g3:first-electric-circuit:bulb}{灯泡}失效：
这正是\omterm{def:g5:series-parallel-circuits:parallel}{并联}规律所体现的独立性。
\textbf{10.} 在所有\omterm{def:g5:series-parallel-circuits:parallel}{支路}的上游发生\omterm{def:g6:circuits-and-safety:short}{短路}。\omterm{def:g7:circuits-series-parallel:current}{电流}
放弃所有正常工作的通路，转而走那条毫无阻力的
裸导线捷径：灯和喇叭得不到\omterm{def:g7:circuits-series-parallel:current}{电流}，而捷径导线中
则涌过一股猛烈发热的\omterm{def:g7:circuits-series-parallel:current}{电流}——所有用电任务同时
失效，直至把该导线拆除。
\textbf{11.} 让大灯在第二组\omterm{def:g3:first-electric-circuit:battery}{电池}上拥有自己独立的回路：
第二组\omterm{def:g3:first-electric-circuit:battery}{电池}、它自己的\omterm{ex:g3:first-electric-circuit:switch}{开关}、大灯——与总配电盘
不共用任何导线。若第一组\omterm{def:g3:first-electric-circuit:battery}{电池}上发生\omterm{def:g6:circuits-and-safety:short}{短路}，
则除船只的灯\emph{以外}的一切都会熄灭。
\textbf{12.} 例如：“\omterm{def:g5:series-parallel-circuits:parallel}{并联}规律让每项任务在今夜保持独立。
今夜再次驳倒了\omterm{def:g3:first-electric-circuit:bulb}{灯泡}会耗尽\omterm{def:g7:circuits-series-parallel:current}{电流}的看法：即便
\omterm{def:g3:first-electric-circuit:battery}{电池}组耗尽，\omterm{def:g7:circuits-series-parallel:current}{电流}仍完整地回归。
一根裸导线使\omterm{def:g3:first-electric-circuit:battery}{电池}组\omterm{def:g6:circuits-and-safety:short}{短路}，让我们经历了
最糟糕的一个\omterm{ex:g2:measuring-time:clock}{小时}：所有通路都空了，而捷径仍在燃烧。”
\end{solution}
